\newpage
\section{Estructura del Código y Manual de Usuario}
\label{sec:manual}

\subsection{Estructura del Proyecto}

El código de la práctica se organiza en el directorio \texttt{code/} con estructura
modular:

\begin{itemize}[noitemsep]
    \item \texttt{common/}: interfaces base del framework (\texttt{MH}, \texttt{Problem}, \texttt{Solution}, utilidades).
    \item \texttt{inc/}: cabeceras de algoritmos y utilidades del problema.
    \item \texttt{src/}: implementaciones C++.
    \item \texttt{datos\_portfolio\_2526/}: ficheros CSV de IBEX 35, S\&P 100 y S\&P 500.
    \item \texttt{config.cfg}: parametrización de experimentos sin recompilar.
    \item \texttt{main.cpp}: orquestación global, ejecución y generación de tablas/CSV.
    \item \texttt{genera\_boxplots.py}: script para figuras de distribución\footnote{Cabe destacar que se ha mejorado/ajustado para este caso.}.
\end{itemize}

\subsection{Módulos Principales Implementados}

\begin{itemize}[noitemsep]
    \item Baselines: \texttt{greedy.cpp}, \texttt{randomsearch.h}, \texttt{localsearch.cpp}.
    \item Enfriamiento Simulado: \texttt{simulated\_annealing.h/.cpp}.
    \item Búsqueda Multiarranque: \texttt{bmb.h/.cpp}.
    \item Búsqueda Local Reiterada: \texttt{ils.h/.cpp}.
    \item Hibridación ILS-ES: \texttt{ils\_es.h/.cpp}.
\end{itemize}

\subsection{Manual de Usuario: Compilación y Ejecución}

El proceso de compilación se ha automatizado completamente mediante \texttt{CMake} y \texttt{Make}. Para compilar el código, el usuario debe situarse en el directorio \texttt{code/} y ejecutar:

\begin{lstlisting}[language=bash, frame=single]
cmake .
make
\end{lstlisting}

Este proceso genera un ejecutable llamado \texttt{main}.

\textbf{Ejecución estándar:} Al iniciarse, el programa lee automáticamente el archivo \texttt{config.cfg}, por lo que no es necesario recompilar para cambiar hiperparámetros.

\begin{lstlisting}[language=bash, frame=single]
./main
\end{lstlisting}

\textbf{Sobrescritura de semilla por línea de comandos:} La semilla también puede pasarse como argumento, sobrescribiendo la del fichero de configuración:

\begin{lstlisting}[language=bash, frame=single]
./main 42
\end{lstlisting}

\subsection{Configuración de Experimentos sin Recompilación}

El fichero \texttt{config.cfg} permite ajustar todos los parámetros de los algoritmos sin necesidad de recompilar. Entre los nuevos parámetros configurables para la Práctica 3 están:

\begin{itemize}[noitemsep]
    \item \textbf{Enfriamiento Simulado:} \texttt{ES\_MU} y \texttt{ES\_PHI} para la temperatura inicial, \texttt{ES\_TF} para la temperatura final, \texttt{ES\_MAX\_VECINOS} y \texttt{ES\_MAX\_EXITOS} para el bucle interno $L(T)$.
    \item \textbf{Algoritmos de Trayectoria Múltiple:} \texttt{TRAJ\_NUM\_STARTS} (número de arranques), \texttt{TRAJ\_MAX\_EVALS\_BL} (presupuesto por cada BL), \texttt{TRAJ\_MUTATION\_RATE} (porcentaje de activos a mutar en ILS).
    \item \textbf{Flags de algoritmos:} para activar/desactivar bloques obligatorios (AG, AM, DE) y los nuevos de la Práctica 3 (\texttt{EXP\_RUN\_P3}).
\end{itemize}

\subsection{Salida de Resultados y Visualización}

Al finalizar su ejecución, el binario genera un resumen en pantalla que incluye:

\begin{itemize}[noitemsep]
    \item Tabla resumen por mercado con fitness, evaluaciones y tiempos de ejecución de cada algoritmo.
    \item Mejor solución (\textit{fitness}) encontrada durante la ejecución.
\end{itemize}

El programa voltea automáticamente los datos en archivos \texttt{.csv} organizados por mercado, permitiendo el análisis posterior.

Para obtener gráficas de distribución (boxplots) de los resultados, basta con ejecutar:

\begin{lstlisting}[language=bash, frame=single]
python3 genera_boxplots.py
\end{lstlisting}

Esto genera visualizaciones en formato \texttt{.png} que facilitan el análisis comparativo entre algoritmos.
